\documentclass[11pt]{article}

\usepackage[margin=1in]{geometry}
\usepackage[T1]{fontenc}
\usepackage[utf8]{inputenc}
\usepackage{lmodern}
\usepackage{amsmath}
\usepackage{amssymb}
\usepackage{booktabs}
\usepackage{longtable}
\usepackage{array}
\usepackage{enumitem}
\usepackage{xcolor}
\usepackage{hyperref}
\usepackage{listings}

\hypersetup{
  colorlinks=true,
  linkcolor=blue!50!black,
  urlcolor=blue!50!black
}

\lstset{
  basicstyle=\ttfamily\small,
  breaklines=true,
  columns=fullflexible,
  keepspaces=true,
  frame=single,
  showstringspaces=false
}

\title{Noise Estimation and Signal-to-Noise Products in Citlali:\\
Definitions, Correlated Noise, and Proposed Data Products}
\author{}
\date{2026-05-01}

\newcommand{\map}{\boldsymbol{m}}
\newcommand{\noise}{\boldsymbol{n}}
\newcommand{\data}{\boldsymbol{d}}
\newcommand{\sky}{\boldsymbol{s}}
\newcommand{\weight}{\boldsymbol{W}}
\newcommand{\cov}{\boldsymbol{C}}
\newcommand{\mat}[1]{\boldsymbol{#1}}
\newcommand{\trans}{\mathsf{T}}
\newcommand{\inv}{^{-1}}
\newcommand{\expect}[1]{\left\langle #1 \right\rangle}
\newcommand{\var}{\operatorname{Var}}
\newcommand{\diag}{\operatorname{diag}}
\newcommand{\median}{\operatorname{median}}
\newcommand{\std}{\operatorname{std}}
\newcommand{\snr}{\mathrm{S/N}}

\begin{document}
\maketitle
\tableofcontents

\section{Purpose}

Citlali currently produces several objects that are all related to noise:
formal weight maps, jackknife noise realization maps, signal-to-noise maps,
coverage maps, and filtered-map weight normalizations. These objects are
useful, but they do not all answer the same scientific question. In particular,
TolTEC commonly produces oversampled maps with pixel sizes of order
\(1\)--\(2\) arcsec and point-spread functions of order \(5\)--\(10\) arcsec.
In that regime neighboring map pixels are strongly correlated. A per-pixel
variance is not, by itself, the uncertainty of an astrophysical flux estimate
unless the estimator is simply ``read off the map at this position.''

The purpose of this note is to propose a rigorous vocabulary and data model for
noise-related products in Citlali. The goal is to support raw maps, filtered
maps, naive mapmaking, jinc mapmaking, individual observations, and coadded
observations without overloading a single quantity called ``the noise.''

The central principle is:

\begin{quote}
An uncertainty is only well-defined after the measurement estimator has been
specified. A weight map describes a map-value variance; a source-flux
uncertainty must be propagated through the same estimator used to measure the
source flux.
\end{quote}

\section{Science Goal}

For most TolTEC science uses, the final desired quantity is not the value of a
single map pixel. It is the flux, surface brightness, or integrated signal of
an astrophysical source under a specified measurement procedure. The
measurement procedure may be:

\begin{itemize}[leftmargin=1.5em]
\item point-source matched filtering or Wiener filtering,
\item beam-profile fitting,
\item aperture photometry,
\item template fitting,
\item smoothed extended-emission analysis,
\item direct map-value inspection for diagnostic or specialized use.
\end{itemize}

These estimators have different noise properties. A point-source matched filter
uses a known kernel. An extended-source aperture sum uses a user-defined region.
A diffuse-emission map smoothed to a target angular scale uses a smoothing
operator. Since the estimators differ, the corresponding uncertainties differ.

\section{Map Products and Units}

\subsection{Signal Map}

Let \(\map\) denote a Citlali map product. For the present discussion, assume
that the signal map has units
\[
  [\map] = \mathrm{mJy\,beam^{-1}}.
\]
This is the natural unit for point-source amplitudes in a beam-normalized map.
For extended-source analysis, conversion factors such as pixel area divided by
beam area must be included in the estimator.

Raw maps and filtered maps should be treated as different data products:
\[
  \map_{\rm raw}, \qquad \map_{\rm filt}.
\]
Filtering changes both the noise covariance and the response to astrophysical
signals. Therefore a filtered signal-to-noise product must be associated with a
filter response or transfer function.

\subsection{Formal Weight Map}

A formal weight map is an inverse variance predicted by the timestream weights
and mapmaking algebra. Its units should be
\[
  [\weight_{\rm formal}] = (\mathrm{mJy\,beam^{-1}})^{-2}.
\]
The corresponding formal per-pixel variance is
\[
  \sigma^2_{\rm formal}(p) = \frac{1}{W_{\rm formal}(p)}.
\]

This is a map-value variance, not a complete description of correlated map
noise. It is also a model prediction: it is only as accurate as the detector
weights, independence assumptions, filtering model, and mapmaking normalization
used to construct it.

\subsection{Jackknife Noise Realization Maps}

Let \(\noise_j\) be jackknife noise realization map \(j\), for
\(j=1,\ldots,N_{\rm jk}\). The intended units are the same as the signal map:
\[
  [\noise_j] = \mathrm{mJy\,beam^{-1}}.
\]
The jackknife maps are an empirical ensemble of noise-like maps produced by
running sign-flipped or split data through the same mapmaking and filtering
pipeline as the signal map. Ideally the astrophysical signal cancels in the
jackknife construction while instrumental, atmospheric, and processing-induced
noise remains statistically representative of the signal map's noise.

\subsection{Empirical Pixel Variance Map}

The empirical per-pixel variance map is the diagonal of the map covariance
estimated from the jackknife ensemble:
\[
  \widehat{\sigma}^2_{\rm pix}(p)
    =
    \frac{1}{N_{\rm jk} - 1}
    \sum_{j=1}^{N_{\rm jk}}
    \left[
      n_j(p) - \overline{n}(p)
    \right]^2,
\]
where
\[
  \overline{n}(p) =
    \frac{1}{N_{\rm jk}}
    \sum_{j=1}^{N_{\rm jk}} n_j(p).
\]
If the jackknife maps are constructed to have zero mean, the simpler second
moment
\[
  \widehat{\sigma}^2_{\rm pix}(p)
    \simeq
    \frac{1}{N_{\rm jk}}
    \sum_{j=1}^{N_{\rm jk}} n_j^2(p)
\]
may be adequate, but this should be a deliberate convention. The variance units
are
\[
  [\widehat{\sigma}^2_{\rm pix}] = (\mathrm{mJy\,beam^{-1}})^2.
\]

\subsection{Empirical Weight Map}

The fully local empirical inverse-variance weight is
\[
  W_{\rm emp,local}(p) =
    \frac{1}{\widehat{\sigma}^2_{\rm pix}(p)}.
\]
This is direct but can be noisy when \(N_{\rm jk}\) is modest. A more stable
alternative is to preserve the formal weight map's spatial pattern while
calibrating its normalization from the jackknife ensemble:
\[
  r(p) =
    W_{\rm formal}(p)\,\widehat{\sigma}^2_{\rm pix}(p),
\]
\[
  \alpha =
    \left[
      \median_{p\in \Omega} r(p)
    \right]^{-1},
\]
\[
  W_{\rm emp,scalar}(p) =
    \alpha\,W_{\rm formal}(p),
\]
where \(\Omega\) is a valid map support region. This scalar calibration is
dimensionally valid because \(r(p)\) is dimensionless.

Intermediate choices are also possible:

\begin{itemize}[leftmargin=1.5em]
\item smooth \(r(p)\) before applying a local correction,
\item estimate one scale factor per observation, array, or map region,
\item regularize \(W_{\rm emp,local}\) toward \(W_{\rm formal}\),
\item use quantile-clipped or source-masked support regions.
\end{itemize}

\section{Formal Mapmaking Weights}

\subsection{Naive Weighted Map}

For a single output pixel \(p\), suppose samples \(i\) contribute data values
\(d_i\) with inverse variances \(w_i\). In a simple nearest-pixel or
non-kernelized weighted map, the accumulated signal and weight are
\[
  S(p) = \sum_{i\in p} w_i d_i,
\]
\[
  Q(p) = \sum_{i\in p} w_i.
\]
The normalized map value is
\[
  m(p) = \frac{S(p)}{Q(p)}.
\]
If the input samples are independent with
\[
  \var(d_i) = \frac{1}{w_i},
\]
then
\[
  \var[m(p)] =
    \frac{\sum_i w_i^2 \var(d_i)}{Q^2(p)}
    =
    \frac{\sum_i w_i}{Q^2(p)}
    =
    \frac{1}{Q(p)}.
\]
Thus the formal weight is
\[
  W_{\rm formal}(p) = Q(p).
\]

\subsection{Jinc Kernel Mapmaking}

For jinc or other kernelized mapmaking, let \(K_{pi}\) be the mapmaking kernel
value that describes how sample \(i\) contributes to output pixel \(p\). The
accumulators are
\[
  S(p) = \sum_i K_{pi} w_i d_i,
\]
\[
  D(p) = \sum_i K_{pi} w_i,
\]
\[
  Q(p) = \sum_i K_{pi}^2 w_i.
\]
The normalized map value is
\[
  m(p) = \frac{S(p)}{D(p)}.
\]
Assuming independent input samples with \(\var(d_i)=1/w_i\), the variance is
\[
\begin{aligned}
  \var[m(p)]
    &=
    \frac{1}{D^2(p)}
    \sum_i K_{pi}^2 w_i^2 \var(d_i) \\
    &=
    \frac{1}{D^2(p)}
    \sum_i K_{pi}^2 w_i \\
    &=
    \frac{Q(p)}{D^2(p)}.
\end{aligned}
\]
Therefore the formal inverse variance is
\[
  W_{\rm formal}(p) =
    \frac{D^2(p)}{Q(p)}.
\]
This expression is important: the jinc kernel does not need to have unit sum.
The normalization is carried by \(D(p)\), while the noise variance depends on
the squared kernel through \(Q(p)\).

\subsection{Limits of Formal Weights}

The formal weights above are exact only under assumptions that are not exactly
true in real TolTEC reductions:

\begin{itemize}[leftmargin=1.5em]
\item detector samples are not independent after timestream filtering,
\item atmospheric subtraction and common-mode removal induce correlations,
\item flagging and source masking make the effective sample ensemble
non-uniform,
\item mapmaking kernels correlate neighboring pixels,
\item filtering changes the noise covariance and signal response,
\item coaddition combines observations with different empirical noise
properties.
\end{itemize}

Consequently, formal weights are valuable diagnostics and useful starting
points, but they should not be the only uncertainty product used for science
significance.

\section{Map Covariance and Correlated Noise}

\subsection{Full Map Covariance}

Let the map be represented as a vector \(\map\) of length \(N_{\rm pix}\). The
full map noise covariance is
\[
  \cov =
    \expect{
      (\map - \expect{\map})
      (\map - \expect{\map})^\trans
    }.
\]
The diagonal entries,
\[
  C_{pp} = \sigma^2_{\rm pix}(p),
\]
are the per-pixel variances. The off-diagonal entries,
\[
  C_{pq}, \qquad p\ne q,
\]
describe pixel-pixel covariance.

In oversampled maps, the off-diagonal covariance is not a small technicality.
It is expected from:

\begin{itemize}[leftmargin=1.5em]
\item the telescope beam,
\item jinc or other mapmaking kernels,
\item timestream filtering,
\item atmospheric/common-mode subtraction,
\item low-pass, matched, or Wiener filtering,
\item coadding with spatially varying coverage.
\end{itemize}

The empirical covariance estimated from jackknife maps is
\[
  \widehat{C}_{pq} =
    \frac{1}{N_{\rm jk}-1}
    \sum_{j=1}^{N_{\rm jk}}
    \left[n_j(p)-\overline{n}(p)\right]
    \left[n_j(q)-\overline{n}(q)\right].
\]
Storing this full matrix is usually impossible for real maps because it scales
as \(N_{\rm pix}^2\). Citlali therefore needs estimator-based covariance
propagation rather than full covariance storage.

\subsection{Linear Estimators}

Many source measurements can be written as linear estimators:
\[
  \widehat{F} =
    \sum_p a_p m(p)
    =
    \mat{a}^\trans \map,
\]
where \(\mat{a}\) is the estimator weight vector. Examples include aperture
photometry, matched filtering evaluated at a position, and smoothed map values.

The variance of this estimator is
\[
  \sigma_F^2 =
    \mat{a}^\trans \cov \mat{a}.
\]
Expanded in pixel notation,
\[
  \sigma_F^2 =
    \sum_p a_p^2 C_{pp}
    +
    2\sum_{p<q} a_p a_q C_{pq}.
\]
The second term is the correlated-noise contribution. If one uses only the
variance map, one implicitly approximates
\[
  \sigma_F^2 \approx \sum_p a_p^2 C_{pp},
\]
which is generally wrong for oversampled maps.

\subsection{Jackknife Propagation Through Estimators}

The practical way to include correlated noise without storing \(\cov\) is to
apply the same estimator to every jackknife realization:
\[
  F_j =
    \mat{a}^\trans \noise_j.
\]
Then
\[
  \widehat{\sigma}_F^2 =
    \frac{1}{N_{\rm jk}-1}
    \sum_{j=1}^{N_{\rm jk}}
    \left(F_j - \overline{F}\right)^2.
\]
This is equivalent to estimating \(\mat{a}^\trans \cov \mat{a}\) from the
jackknife ensemble. It automatically includes pixel-pixel covariance, the jinc
kernel, filtering, coadd weighting, edge effects, and coverage variation to the
extent that these effects are present in the jackknife maps.

This leads to the main operational rule:

\begin{quote}
For a science uncertainty, run the science estimator on the jackknife ensemble
and use the scatter of the resulting measurements.
\end{quote}

\section{Signal-to-Noise Maps}

\subsection{Pixel Signal-to-Noise}

A pixel signal-to-noise map is
\[
  \snr_{\rm pix}(p) =
    \frac{m(p)}{\widehat{\sigma}_{\rm pix}(p)}
    =
    m(p)\sqrt{W_{\rm emp}(p)}.
\]
This is dimensionless. It answers: ``How significant is the map value at this
pixel relative to the empirical per-pixel RMS?''

It does not, in general, answer: ``How significant is a source flux
measurement?'' That depends on the source estimator.

\subsection{Point-Source Signal-to-Noise}

For point sources, the source shape is approximately known: the effective beam
or filtered point-source response. Let \(\mat{t}_{x_0}\) be a point-source
template centered at position \(x_0\). If the map covariance were known, the
minimum-variance linear amplitude estimator would be
\[
  \widehat{F}(x_0) =
    \frac{
      \mat{t}_{x_0}^\trans \cov\inv \map
    }{
      \mat{t}_{x_0}^\trans \cov\inv \mat{t}_{x_0}
    },
\]
with variance
\[
  \sigma_F^2(x_0) =
    \left(
      \mat{t}_{x_0}^\trans \cov\inv \mat{t}_{x_0}
    \right)^{-1}.
\]
In practice Citlali may not have or store \(\cov\inv\). A matched filter or
Wiener filter approximates this operation under a chosen noise model. The
empirical uncertainty can still be obtained by applying the same point-source
filter and response normalization to every jackknife map:
\[
  F_j(x_0) =
    {\cal E}_{\rm ps}[\noise_j](x_0),
\]
\[
  \widehat{\sigma}_{F,\rm ps}^2(x_0) =
    \var_j \left[ F_j(x_0) \right].
\]
This point-source matched/Wiener-filter case is the special case where Citlali
is closest to treating the source noise estimate correctly as part of the map
product itself. The source template is specified by the beam or effective
kernel, the filter is designed for that template, and the filtered jackknife
maps can be propagated through the same operation as the signal map. Under
those conditions, the relevant noise is no longer just the diagonal per-pixel
map variance; it is the variance of the filtered point-source amplitude
estimator. This statement depends on the filter response normalization and the
jackknife noise normalization being applied consistently.

The point-source signal-to-noise is then
\[
  \snr_{\rm ps}(x_0) =
    \frac{
      {\cal E}_{\rm ps}[\map](x_0)
    }{
      \widehat{\sigma}_{F,\rm ps}(x_0)
    }.
\]

\subsection{Extended-Source Signal-to-Noise}

There is no universal extended-source flux uncertainty because there is no
universal extended-source shape. An extended-source measurement must specify an
estimator. For an aperture or mask \(R\), a common integrated-flux estimator is
\[
  \widehat{F}_R =
    \sum_{p\in R}
    m(p)\,
    \frac{\Omega_{\rm pix}}{\Omega_{\rm beam}},
\]
where \(\Omega_{\rm pix}\) is the pixel solid angle and \(\Omega_{\rm beam}\)
is the beam solid angle. The jackknife propagated uncertainty is
\[
  \widehat{\sigma}_{F,R}^2 =
    \var_j
    \left[
      \sum_{p\in R}
      n_j(p)\,
      \frac{\Omega_{\rm pix}}{\Omega_{\rm beam}}
    \right].
\]
This includes pixel-pixel covariance inside the aperture. A diagonal-only
approximation would be
\[
  \sigma^2_{F,R,{\rm diag}} =
    \sum_{p\in R}
    \widehat{\sigma}^2_{\rm pix}(p)
    \left(
      \frac{\Omega_{\rm pix}}{\Omega_{\rm beam}}
    \right)^2,
\]
but this is not expected to be accurate when the map is oversampled.

For smoothed diffuse-emission products, define a smoothing operator
\({\cal S}_L\), such as convolution to a target scale \(L\):
\[
  \map_L = {\cal S}_L[\map].
\]
The corresponding noise map ensemble is
\[
  \noise_{j,L} = {\cal S}_L[\noise_j],
\]
and the smoothed-map variance is
\[
  \widehat{\sigma}^2_L(p) =
    \var_j \left[ n_{j,L}(p) \right].
\]
The scale \(L\) is part of the definition of the product.

\section{Raw and Filtered Maps}

\subsection{Raw Maps}

Raw maps should expose both formal and empirical noise products:
\[
  \map_{\rm raw},
  \quad
  W_{{\rm raw},{\rm formal}},
  \quad
  \{\noise_{{\rm raw},j}\},
  \quad
  \widehat{\sigma}^2_{{\rm raw},{\rm pix}},
  \quad
  W_{{\rm raw},{\rm emp}},
  \quad
  \snr_{{\rm raw},{\rm pix}}.
\]
The formal raw weights are useful for coadd weighting and diagnostics, but the
empirical calibration determines whether the raw pixel signal-to-noise has unit
noise width.

If the formal raw weights overestimate the empirical inverse variance by an
approximately global factor, one can define
\[
  \alpha_{\rm raw}
    =
    \left[
      \median_{p\in\Omega}
      W_{{\rm raw},{\rm formal}}(p)
      \widehat{\sigma}^2_{{\rm raw},{\rm pix}}(p)
    \right]^{-1},
\]
and
\[
  W_{{\rm raw},{\rm emp}}(p)
    =
    \alpha_{\rm raw}
    W_{{\rm raw},{\rm formal}}(p).
\]

\subsection{Filtered Maps}

A filtered map has a different covariance and a different astrophysical
response:
\[
  \map_{\rm filt} = {\cal F}[\map_{\rm raw}],
\]
\[
  \noise_{{\rm filt},j} = {\cal F}[\noise_{{\rm raw},j}].
\]
The filtered-map empirical variance is
\[
  \widehat{\sigma}^2_{{\rm filt},{\rm pix}}(p)
    =
    \var_j[n_{{\rm filt},j}(p)].
\]

If the filter is intended for point-source work, the filtered signal map should
also carry a point-source response normalization. Let \(R_{\rm ps}(p)\) be the
response of the filtered product to a unit-flux point source at position \(p\).
Then a point-source flux map is schematically
\[
  F_{\rm ps}(p) =
    \frac{\map_{\rm filt}(p)}{R_{\rm ps}(p)}.
\]
The noise realizations should be normalized in the same way:
\[
  F_{{\rm ps},j}(p) =
    \frac{n_{{\rm filt},j}(p)}{R_{\rm ps}(p)}.
\]
The point-source uncertainty is
\[
  \widehat{\sigma}^2_{{\rm ps}}(p)
    =
    \var_j[F_{{\rm ps},j}(p)].
\]

\section{Individual Observations and Coadds}

\subsection{Individual Observation Products}

For each observation \(o\), Citlali can define:
\[
  \map_o,
  \quad
  W_{o,{\rm formal}},
  \quad
  \{\noise_{o,j}\},
  \quad
  W_{o,{\rm emp}},
  \quad
  \widehat{\sigma}^2_{o,{\rm pix}}.
\]
These products should be computed in the same mapmaking mode as the signal map:
naive, jinc, or another future mapmaker.

\subsection{Coadded Signal Maps}

Suppose observation maps are coadded with weights \(A_o(p)\). A common
inverse-variance coadd is
\[
  \map_{\rm coadd}(p)
    =
    \frac{
      \sum_o A_o(p)\,\map_o(p)
    }{
      \sum_o A_o(p)
    }.
\]
If the observation maps are empirically calibrated, a natural choice is
\[
  A_o(p) = W_{o,{\rm emp}}(p).
\]
If only formal weights are available, then the coadd may inherit the formal
weight calibration errors of the inputs.

\subsection{Coadded Noise Realizations}

The coadded jackknife ensemble should be constructed by coadding the observation
noise maps with the same coadd rule:
\[
  \noise_{{\rm coadd},j}(p)
    =
    \frac{
      \sum_o A_o(p)\,\noise_{o,j}(p)
    }{
      \sum_o A_o(p)
    }.
\]
Then
\[
  \widehat{\sigma}^2_{{\rm coadd},{\rm pix}}(p)
    =
    \var_j[n_{{\rm coadd},j}(p)].
\]
This final empirical coadd variance is the most robust calibration of the
coadded map-value noise because it includes the actual coadd operation.

\subsection{Indexing Jackknife Realizations Across Observations}

There are several possible ways to combine observation-level jackknife
realizations:

\begin{enumerate}[leftmargin=1.5em]
\item Use common realization index \(j\) across observations and form
\(\noise_{{\rm coadd},j}\) from all \(\noise_{o,j}\).
\item Randomly assign or resample observation-level jackknife signs to construct
a larger coadd noise ensemble.
\item Generate coadd-level jackknife splits directly from the combined data.
\end{enumerate}

The first option is simple and deterministic, but it assumes that realization
index \(j\) has compatible meaning across observations. The second can improve
sampling but requires clear provenance. The third is conceptually clean but may
be harder to implement in Citlali's current reduction structure.

The important requirement is that the coadd noise ensemble must represent the
same coadd operation used for the signal map.

\section{Memory and Disk Management}

\subsection{Problem}

Writing every possible noise-related product can be expensive. If a reduction
has \(N_{\rm jk}\) noise maps, \(N_{\rm band}\) arrays, several filtered
products, and both observation-level and coadd-level outputs, the disk volume
can grow by a large multiplicative factor. In-memory storage can be similarly
problematic.

Citlali should therefore distinguish:

\begin{itemize}[leftmargin=1.5em]
\item full noise realizations,
\item streaming accumulators,
\item compact derived summaries,
\item estimator-specific uncertainty products.
\end{itemize}

\subsection{Streaming Accumulators}

For pixel variance, Citlali does not need to retain all jackknife maps in memory
if it only needs the first two moments. For each pixel:
\[
  A_1(p) = \sum_{j=1}^{N_{\rm jk}} n_j(p),
\]
\[
  A_2(p) = \sum_{j=1}^{N_{\rm jk}} n_j^2(p).
\]
Then
\[
  \widehat{\sigma}^2_{\rm pix}(p)
    =
    \frac{
      A_2(p) - A_1^2(p)/N_{\rm jk}
    }{
      N_{\rm jk}-1
    }.
\]
This requires only two map-sized accumulators per product.

For a scalar-calibrated weight normalization, Citlali additionally needs a
robust statistic of
\[
  r(p)=W_{\rm formal}(p)\widehat{\sigma}^2_{\rm pix}(p)
\]
over the valid support. This can be computed after the variance map is formed.

\subsection{Estimator Accumulators}

For any estimator \({\cal E}\), Citlali can stream jackknife realizations
through the estimator:
\[
  F_j = {\cal E}[\noise_j].
\]
If \(F_j\) is a scalar, such as an aperture flux for one region, retain only
\[
  B_1 = \sum_j F_j,
  \qquad
  B_2 = \sum_j F_j^2.
\]
If \(F_j(p)\) is a map, such as a point-source flux map, retain map-sized
accumulators
\[
  B_1(p) = \sum_j F_j(p),
  \qquad
  B_2(p) = \sum_j F_j^2(p).
\]
The estimator variance follows from the same two-moment formula.

\subsection{Recommended Write Policy}

The default on-disk products should be compact. A possible default set is:

\begin{itemize}[leftmargin=1.5em]
\item signal map,
\item formal weight map,
\item empirical weight map or empirical scalar normalization,
\item coverage map,
\item pixel signal-to-noise map,
\item point-source signal-to-noise map when a point-source filter is requested,
\item compact noise summary metadata or table.
\end{itemize}

Full jackknife noise realization maps should be optional:
\[
  \texttt{write\_noise\_realizations = true/false}.
\]
They are extremely useful for diagnostics and downstream custom analyses, but
they should not be required for every routine reduction if Citlali can stream
the necessary uncertainty summaries.

\section{Proposed Product Taxonomy}

\begin{longtable}{p{0.22\linewidth} p{0.19\linewidth} p{0.50\linewidth}}
\toprule
Product & Units & Meaning \\
\midrule
\endhead
\(\map\) & \(\mathrm{mJy\,beam^{-1}}\) & Signal map in the native response of the map product. \\
\(W_{\rm formal}\) & \((\mathrm{mJy\,beam^{-1}})^{-2}\) & Inverse variance predicted by timestream weights and mapmaking algebra. \\
\(\noise_j\) & \(\mathrm{mJy\,beam^{-1}}\) & Jackknife noise realization map produced through the same pipeline path as the signal product. \\
\(\widehat{\sigma}^2_{\rm pix}\) & \((\mathrm{mJy\,beam^{-1}})^2\) & Empirical per-pixel map-value variance from the jackknife ensemble. \\
\(W_{\rm emp}\) & \((\mathrm{mJy\,beam^{-1}})^{-2}\) & Empirically calibrated inverse variance, either local or scalar-renormalized. \\
\(\snr_{\rm pix}\) & dimensionless & Pixel signal-to-noise, \(m(p)\sqrt{W_{\rm emp}(p)}\). \\
\(R_{\rm ps}\) & map units per flux unit & Point-source response of a filtered product. \\
\(F_{\rm ps}\) & mJy & Point-source flux estimate after response normalization. \\
\(\sigma_{F,\rm ps}\) & mJy & Point-source flux uncertainty from jackknife propagation. \\
\(\snr_{\rm ps}\) & dimensionless & Point-source flux signal-to-noise. \\
\(\sigma_{F,R}\) & mJy & Aperture or region-integrated flux uncertainty for specified region \(R\). \\
\bottomrule
\end{longtable}

\section{Suggested Citlali Configuration Model}

The following block is illustrative rather than a final schema:

\begin{lstlisting}
noise:
  enabled: true

  realizations:
    n_noise: 25
    write: false
    write_coadd: false

  pixel_variance:
    enabled: true
    estimator: jackknife
    mean_subtract: true
    support: coverage_bool
    calibration:
      mode: scalar_renorm        # scalar_renorm, local, smoothed_local
      statistic: median
      source_mask: true

  products:
    write_formal_weight: true
    write_empirical_weight: true
    write_pixel_snr: true
    write_noise_summary: true

  estimators:
    point_source:
      enabled: true
      method: wiener             # wiener, matched_filter, beam_fit
      write_flux_map: true
      write_uncertainty_map: true
      write_snr_map: true

    aperture:
      enabled: false
      masks: []
      write_catalog: true

    smoothed_surface_brightness:
      enabled: false
      fwhm_arcsec: []
      write_variance_maps: true
\end{lstlisting}

\section{Diagnostics}

Citlali should report noise diagnostics that make weight failures visible
without requiring users to inspect every map manually. Useful diagnostics
include:

\begin{itemize}[leftmargin=1.5em]
\item \(\median_p\left[
  W_{\rm formal}(p)\widehat{\sigma}^2_{\rm pix}(p)
\right]\),
\item robust scatter of \(W_{\rm formal}\widehat{\sigma}^2_{\rm pix}\),
\item histogram width of \(n_j(p)\sqrt{W_{\rm formal}(p)}\),
\item histogram width of \(n_j(p)\sqrt{W_{\rm emp}(p)}\),
\item histogram width of \(m(p)\sqrt{W_{\rm emp}(p)}\),
\item number of valid pixels in the support region,
\item fraction of pixels with non-finite or zero variance,
\item point-source noise width from filtered jackknife maps,
\item comparison of formal and empirical coadd weights.
\end{itemize}

The distinction between signal-map histograms and jackknife-noise histograms is
important. A broad signal-map \(\snr\) distribution may indicate real sky
structure, residual background, source contamination, or edge artifacts. A
broad jackknife-noise \(\snr\) distribution more directly indicates a noise
normalization problem.

\section{Recommended Development Phases}

\subsection{Phase 1: Definitions and Metadata}

Adopt explicit names for formal weights, empirical weights, pixel variance,
pixel S/N, point-source S/N, and noise realizations. Add metadata describing:

\begin{itemize}[leftmargin=1.5em]
\item whether a weight map is formal or empirical,
\item how empirical calibration was computed,
\item which support region was used,
\item whether the jackknife variance used mean subtraction,
\item which estimator defines an S/N product,
\item which filter response normalization was applied.
\end{itemize}

\subsection{Phase 2: Raw and Filtered Pixel Variance}

Implement streaming two-moment accumulators for raw and filtered jackknife maps.
Use them to produce:
\[
  \widehat{\sigma}^2_{\rm pix},
  \quad
  W_{\rm emp},
  \quad
  \snr_{\rm pix}.
\]
Keep full noise-realization writing optional.

\subsection{Phase 3: Point-Source Estimator Products}

Define the point-source estimator used by filtered maps. Apply the same
estimator and response normalization to signal and jackknife maps. Produce:
\[
  F_{\rm ps},
  \quad
  \sigma_{F,\rm ps},
  \quad
  \snr_{\rm ps}.
\]
This is likely the most important science-facing uncertainty product for point
source catalogs.

\subsection{Phase 4: Coadd Noise Propagation}

Construct coadd-level jackknife products by applying the same coadd weights and
support masks used for the signal coadd. Calibrate final coadd empirical
weights from the coadd noise ensemble.

\subsection{Phase 5: Extended-Source Estimator Interface}

Add an estimator interface for aperture, template, or smoothed-map uncertainty
products. Extended-source products should require explicit estimator
definitions rather than relying on a generic source-flux uncertainty.

\section{Open Questions}

\begin{enumerate}[leftmargin=1.5em]
\item Should empirical pixel variance use the jackknife second moment or
sample variance with mean subtraction by default?
\item Should raw empirical weights default to scalar renormalization or local
variance maps?
\item What support region should be used for scalar calibration: coverage cut,
edge-guarded region, source-masked region, or high-coverage core?
\item Should individual-observation empirical weights be used in coadding, or
should the coadd be calibrated only after formal coaddition?
\item How should jackknife realization indices be matched or resampled across
observations in coadds?
\item What is the canonical point-source response normalization for each
filtered product?
\item Which products should be default on disk, and which should be optional
debug/provenance products?
\item How should Citlali expose extended-source estimators without forcing a
single definition of extended-source flux?
\end{enumerate}

\section{Summary}

Citlali should treat noise estimation as an estimator-dependent problem. The
formal weight map is a useful mapmaker prediction, but jackknife realizations
are the empirical noise ensemble. The diagonal variance of that ensemble gives
the uncertainty of map values, while source-flux uncertainties require applying
the source estimator to the same ensemble.

For oversampled TolTEC maps, correlated noise is unavoidable. The mathematically
correct variance of a linear estimator is
\[
  \sigma_F^2 = \mat{a}^\trans \cov \mat{a},
\]
not merely a sum over per-pixel variances. Since storing \(\cov\) is generally
impractical, the practical solution is to propagate jackknife noise
realizations through the same estimators used on the signal map. This approach
handles naive and jinc mapmaking, raw and filtered maps, individual
observations, and coadds, while allowing Citlali to control memory and disk use
through streaming accumulators and optional realization writing.

\end{document}
